\section[Examples]{\added{Examples}}
\label{sec:examples}

\Added{This section runs the interpretation at particular semirings, on particular programs, to show the
analyses that result. Each example is mechanised in the accompanying Agda development.}

\subsection[Perturbation Bounds]{\added{Perturbation Bounds}}
\label{sec:examples:perturbation-bounds}

\Added{The first two instances demonstrate the perturbation-bound reading of
\secref{approx-as-tangents:richer-semirings}: the analysis reports how much each output can change, rather
than whether it can. Sections \ref{sec:examples:absolute-bounds} and \ref{sec:examples:relative-bounds}
give the two presentations, absolute and relative, run on the database $\mathit{db} = [(\mathsf{a}, 3),
(\mathsf{b}, 1), (\mathsf{a}, -3)]$ of \secref{introduction:galois-slicing}. In both, the approximation
object for a number is the free semimodule $S^2$, one bound for each direction of change.}

\subsubsection[Absolute bounds]{\added{Absolute bounds}}
\label{sec:examples:absolute-bounds}

\Added{At the min-plus semiring, a pair of scalars $(d, u)$ records that a value may fall by at most $d$ and
rise by at most $u$, so that a number $x$ is known to lie in the interval $[x - d, x + u]$; the pair
$(\infty, \infty)$ carries no information.}

\Added{Multiplication's derivative carries no information at this semiring, since no translation bounds a
scaling (\secref{approx-as-tangents:richer-semirings}), so we run the additive part of the running example:
the selection-and-sum query $\mathrm{query}\,l\,\mathit{db} = \mathrm{sum}\,[ q \mid (l', q) \leftarrow
\mathit{db}, l \equiv l' ]$, with $\mathrm{query}\,\mathsf{a}\,\mathit{db}
= 0$. Addition's derivative entries are translations by $0$, and selection's entries are the semiring's $0$
and $1$, as before. The backward derivative at this run is a linear map}
\begin{AddedBlock}
\begin{displaymath}
  \partial(\mathrm{query}\,l\,\mathit{db})_r : S^2 \multimap S^2 \times S^2 \times S^2
\end{displaymath}
\end{AddedBlock}
\noindent\Added{with one factor for each number in the database (the frozen label positions and the list terminator
contribute trivial factors, omitted here). Because the input is a list, the type of the derivative depends on
the shape of the input.}

\Added{Backwards, asking for the output to lie within $[-\frac{1}{10}, \frac{1}{10}]$, that is, bounds
$(\frac{1}{10}, \frac{1}{10})$, returns the same bounds at both $\mathsf{a}$-numbers and $(\infty,
\infty)$ at the $\mathsf{b}$-number: changing either $\mathsf{a}$-number alone by at most $\frac{1}{10}$
keeps the output within its interval, while the $\mathsf{b}$-number is unconstrained. Forwards, if the first
$\mathsf{a}$-number is known to within bounds $(\frac{1}{2}, 0)$ and the second to within $(\frac{1}{5},
\frac{1}{2})$, the output is known to within $(\frac{1}{5}, 0)$: bounds for the same position combine by
$\min$, reflecting the reading in which one input changes at a time.}

\subsubsection[Relative bounds]{\added{Relative bounds}}
\label{sec:examples:relative-bounds}

\Added{In the second presentation a scalar bounds the relative rather than the absolute change of a value,
and the two readings trade places over the arithmetic operations: relative bounds pass through multiplication
unchanged, but are amplified by addition. Because the absolute presentation could say nothing about
multiplication, it exercised only the additive part of the running example. With relative bounds we can run
the full weighted-sum query $\mathrm{total}$ of \secref{introduction:galois-slicing}, prices included, on
the same database.}

\Added{We instantiate at the min-times semiring. The derivative of multiplication now has entries $1$ on both
arguments, since a bound on the relative error of either factor is a bound on the relative error of the
product. The derivative of addition at values $(x, y)$ has entry $|x| / |x + y|$ on its first argument and
symmetrically on its second: the factor by which the relative error of a summand is amplified in the sum. When
$x + y = 0$ the factor is infinite, and the entry is the semiring zero $\infty$.}

\Added{In the run of $\mathrm{total}\,\mathsf{b}$, multiplication behaves well: if the price of
$\mathsf{b}$ is known to within $10\%$, the output is too, and likewise for the selected quantity;
backwards, an output bound of $10\%$ requires the used inputs to within $10\%$ and constrains nothing else.
The run of $\mathrm{total}\,\mathsf{a}$ meets the cancellation: its two contributions, $2 \cdot 3$ and $2
\cdot (-3)$, sum to zero, the amplification factor of the final addition is infinite, and no relative bound
on either $\mathsf{a}$-quantity, however tight, yields a bound on the output. The cancellation that the
Boolean instance reported as a spurious dependency in \secref{introduction:galois-slicing} is here reported
as an infinite condition number, in the sense of numerical analysis~\cite{high:ASNA2}.}

\Added{Neither reading of perturbation bounds is better than the other. Absolute bounds propagate
accurately through addition but carry no information through multiplication, while relative bounds propagate
accurately through multiplication but are amplified by addition, and lost entirely where contributions
cancel. The choice between the two presentations depends on what the analysis is meant to observe.}

\subsection[Linked Inputs and Outputs]{\added{Linked Inputs and Outputs}}
\label{sec:examples:linked}

\Added{When outputs overlap in their dependencies, the two derivatives compose into a further analysis: a
map sending a selection of outputs to the outputs related to it by a shared dependency. The
\emph{moving average} is the standard example of such overlap. With window two and four inputs,}
\begin{AddedBlock}
\begin{displaymath}
  \mathrm{mavg}\,(x_1, x_2, x_3, x_4) =
  (\tfrac{1}{2} (x_1 + x_2),\; \tfrac{1}{2} (x_2 + x_3),\; \tfrac{1}{2} (x_3 + x_4))
\end{displaymath}
\end{AddedBlock}
\noindent\Added{Each adjacent pair of outputs shares an input, and the
first and last outputs share none. Running over the Booleans at the input $(1, 2, 4, 8)$, the Jacobian and
the composite of the backward and forward derivatives are}
\begin{AddedBlock}
\begin{displaymath}
  \partial_{\Two}(\mathrm{mavg}) =
  \begin{pmatrix} 1 & 1 & 0 & 0 \\ 0 & 1 & 1 & 0 \\ 0 & 0 & 1 & 1 \end{pmatrix}
  \qquad
  \partial_{\Two}(\mathrm{mavg}) \comp \transpose{\partial_{\Two}(\mathrm{mavg})} =
  \begin{pmatrix} 1 & 1 & 0 \\ 1 & 1 & 1 \\ 0 & 1 & 1 \end{pmatrix}
\end{displaymath}
\end{AddedBlock}
\Added{The composite relates each output to the outputs it shares an input with: selecting the first output
returns the first and second, and the third stays $\bot$ because the two windows are disjoint. This is the
\emph{cognacy} analysis of linked visualisations~\cite{perera22,bond25,psallidas18}, computed here as a
composite of the two derivatives of one run.}

\subsection[Counting]{\added{Counting}}
\label{sec:examples:counting}

\Added{The counting semiring $(\mathbbm{N}, +, 0, \cdot, 1)$ of
\secref{approx-as-tangents:richer-semirings} refines dependency tracking from whether an input is used to
how many times. Every argument of every operation is given the coefficient $1$, so a Jacobian entry counts
the paths from an input position to an output position.}

\Added{On the weighted-sum query of \secref{introduction:galois-slicing}, the backward derivative of the
output returns a count for every position: $1$ at each selected quantity, $0$ at the unselected row and the
other price, and $2$ at the price of the queried label, which is used once in each selected row. The count
of $2$ survives the cancellation that zeroes the rational derivative, because the counting semiring is
positive and counts cannot cancel. The three instances thus answer three different questions at the same
position: the rational Jacobian reports sensitivity ($0$), the Boolean Jacobian reports possible dependence
($1$), and the counting Jacobian reports usage ($2$). Zero-testing sends the count $2$ to the Boolean $1$
exactly, as \secref{approx-as-tangents:richer-semirings} leads us to expect of a positive semiring, whereas
the same specialisation applied to the rational entry loses the dependency.}
